A factory has $n$ machines which can be used to make products. Your goal is to make a total of $t$ products.
For each machine, you know the number of seconds it needs to make a single product. The machines can work simultaneously, and you can freely decide their schedule.
What is the shortest time needed to make $t$ products?

\vspace{0.3cm}\noindent\textbf{Input}

The first input line has two integers $n$ and $t$: the number of machines and products.
The next line has $n$ integers $k_1,k_2,\dots,k_n$: the time needed to make a product using each machine.

\vspace{0.3cm}\noindent\textbf{Output}

Print one integer: the minimum time needed to make $t$ products.

\vspace{0.3cm}\noindent\textbf{Constraints}

\textbullet\ $1 \le n \le 2 \cdot 10^5$
\textbullet\ $1 \le t \le 10^9$
\textbullet\ $1 \le k_i \le 10^9$